\begingroup
\newcommand{\TransBracket}[1]{\boldsymbol{[}#1\boldsymbol{]}}
\newcommand{\IndexSurd}[1]{{}^{#1}\!\surd}
\ArithmeticHead{302}{L}{ANALYSIS. NACHLASS.}
\NumberedItem{13}
\par\indent Zu den Hauptsätzen in dieser Theorie gehört folgendes Theorem.
\par\indent Bezeichnet man
\[\begin{aligned}
&\Bigl(e^{-\frac{P\alpha}{24}}-e^{-\frac{25P\alpha}{24}}-e^{-\frac{49P\alpha}{24}}\\
&\qquad+e^{-\frac{121P\alpha}{24}}+e^{-\frac{169P\alpha}{24}}-{.}{.}{.}\Bigr)\IndexSurd{4}\alpha\ \text{durch }\varphi\alpha,
\end{aligned}\]
\noindent so ist
\[
\varphi\alpha=\varphi\frac1\alpha\mathpunct{\raisebox{.5ex}{.}}
\]
\noindent Diese Function ist ein Maximum für $\alpha=1$, wo ihr Werth
\[
=0{,}7682255,\quad\log\,{.}{.}{.}\ 9{,}8854887;
\]
\noindent das Quadrat [ist $=$] $0{,}5901704$, $\surd\frac12$ dadurch dividirt
\[
=1{,}19814=\frac\pi\varpi\ \text{[*)]}.
\]
\Needspace{8\baselineskip}
\par\smallskip\noindent\rule{.25\linewidth}{.4pt}\par
{\small
\par\indent [*)] Das obige Maximum [ist] also
\[
=\frac1{\IndexSurd{4}2\surd1{,}19814{.}{.}{.}}\mathpunct{\raisebox{.5ex}{.}}
\]
\noindent Hier das numerische:
\begin{center}
\setlength{\tabcolsep}{5pt}
\renewcommand{\arraystretch}{1.12}
\begin{tabular}{@{}l r|r|r|}
$\log P\ { .}\,{.}\,{.}\,{.}\,{.}$&$0{,}798179868358$&$1$&$9{,}8863019705$\\
$\log M\ { .}\,{.}\,{.}\,{.}\,{.}$&$9{,}6377843114$&$25$&$7{,}1575492622$\\
$\mathrm{C}.\log24\ { .}\,{.}\,{.}$&$8{,}619788758288$&$49$&$4{,}4287965540$\\\cline{2-2}
&$9{,}055752938046$&$121$&$6{,}2425384293$\\
\multicolumn{1}{r}{Zahl}&$0{,}11369802951$&\multicolumn{2}{c}{}\\
\end{tabular}
\par\medskip
\begin{tabular}{@{}r|r|r r@{}}
\multicolumn{1}{r}{}&\multicolumn{1}{c}{doppelt}&\multicolumn{2}{c}{halb}\\
$0{,}7696654$&$0{,}5923848$&&$0{,}8773058$\\
$14373\;1$&$21$&$-$&$379118$\\
$26\;8$&&$-$&$16383$\\\cline{1-1}
$0{,}7682254$&&$+$&$1$\\\cline{3-4}
\multicolumn{1}{r}{}&\multicolumn{1}{r}{$9{,}7726024$}&&$0{,}8377558$\\
\multicolumn{1}{r}{$9{,}8854887$}&\multicolumn{1}{r}{$0{,}0752575$}&&$0{,}0752575$\\
\multicolumn{1}{r}{}&\multicolumn{1}{r}{$9{,}8478599$}&&$9{,}9231174$\\
\multicolumn{1}{r}{}&\multicolumn{1}{r}{}&&$9{,}8478599$\\
\end{tabular}
\end{center}
\par}
\clearpage
\ArithmeticHead{303}{R}{{\scriptsize ZUR THEORIE DER TRANSSCENDENTEN FUNCTIONEN. I.}}
\SourceLeaf{310}
\noindent Ferner ist
\[
(\varphi2\lambda)^{24}=4(\varphi\lambda)^8(\varphi4\lambda)^8\{(\varphi\lambda)^8+(\varphi4\lambda)^8\}
\]
\[
\varphi\alpha=e^{-\frac\alpha{24}P}\IndexSurd{4}\alpha(1-e^{-\alpha P})(1-e^{-2\alpha P})(1-e^{-3\alpha P}){.}{.}{.}
\]
\noindent oder
\[
\TransBracket{x}=\frac{\displaystyle\varphi\frac{-\log x}{P}\cdot\IndexSurd{24}\frac1x}{\displaystyle\IndexSurd{4}\frac{-\log x}{P}}\mathpunct{\raisebox{.5ex}{.}}
\]
\ClosingRule
\par\indent Es sei
\[\begin{aligned}
&(1+xy)(1+x^3y)(1+x^5y){.}{.}{.}\\
&\qquad\left(1+\frac xy\right)\left(1+\frac{x^3}y\right)\left(1+\frac{x^5}y\right){.}{.}{.}=F(x,y),
\end{aligned}\]
\noindent so ist
\[\begin{aligned}
&F\left(e^{-\frac12\alpha P},e^{P\omega\surd\alpha}\right)e^{-\frac{P}{24\alpha}}\\
&\qquad=F\left(e^{-\frac12\frac P\alpha},e^{iP\omega\surd\frac1\alpha}\right)e^{-\frac{P\alpha}{24}}.\,e^{\frac12P\omega\omega}\ \text{[*)]}.
\end{aligned}\]
\Needspace{5\baselineskip}
\par\smallskip\noindent\rule{.25\linewidth}{.4pt}\par
{\small\par\indent [*)] So ist, $\log x\log x'=\pi\pi$ gesetzt,
\[
x^{\frac1{12}+\omega\omega}F(x,x^{2\omega})=F(x',e^{-iP\omega})x'^{\frac1{12}}.
\]
\par}
\clearpage
\endgroup
